\documentclass[12pt,aspectratio=1610,professionalfont,serif,t]{beamer}
\usepackage[osf,p]{newpxtext}
\usepackage{newpxmath}
\usepackage{fontawesome5}

\usepackage{fancyvrb}
\fvset{%frame=single,
commandchars=\\\{\},
%framesep=1mm,
fontfamily=helvetica,
fontsize=\normalsize
}
%\usepackage[bitstream-charter]{mathdesign}
\usepackage{listings}
\lstset{commentstyle=\color{orange}, keywordstyle=\color{yellow} }


\usetheme{Boadilla}
\useoutertheme{split}
\usecolortheme{albatross}


\setbeamertemplate{blocks}[rounded][shadow=false]
\setbeamersize{text margin left=\dimexpr 1cm+1ex\relax, text margin right=\dimexpr 1cm+1ex\relax}
\addtobeamertemplate{headline}{}{\usebeamerfont{frametitle}\vspace*{1ex}}
\setbeamertemplate{navigation symbols}{}
\setbeamertemplate{page number in head/foot}{}
%\beamerdefaultoverlayspecification{<+->}


\setbeamercolor*{structure}{fg=green!60!black,bg=black}
\setbeamercolor*{normal text}{fg=green!60!black,bg=black}
\setbeamercolor{palette primary}{use={structure,normal text},fg=green,bg=structure.bg}
\setbeamercolor{palette secondary}{use={structure,normal text},fg=structure.fg,bg=structure.bg}
\setbeamercolor{palette tertiary}{use={structure,normal text},fg=structure.fg,bg=structure.bg}
\setbeamercolor{palette quaternary}{use={structure,normal text},fg=white,bg=green!30!black}
\setbeamercolor{palette sidebar quaternary}{use={structure,normal text},fg=white,bg=green!30!black}
\setbeamercolor{title in head/foot}{use={structure,normal text},fg=white,bg=green!30!black}

\setbeamercolor*{example text}{use={structure},fg=green!60!black}
\setbeamercolor*{block body}{use={structure},bg=structure.bg}
\setbeamercolor*{block body alerted}{use={structure},bg=structure.bg}
\setbeamercolor*{block body example}{use={structure},bg=structure.bg}
\setbeamercolor*{block title}{parent=structure,bg=structure.bg}
\setbeamercolor*{block title alerted}{use={structure,alerted text},fg=alerted text.fg,bg=structure.bg}
\setbeamertemplate{navigation symbols}{}
\setbeamertemplate{items}[square]
\setbeamercolor{item projected}{fg=white}
\setbeamercolor*{normal text}{use={structure},fg=white!90!blue,bg=structure.bg}
\setbeamercolor*{separation line}{}
\setbeamercolor*{fine separation line}{}
\setbeamercolor{alerted text}{use={structure},fg=structure.fg}

\usepackage[italian]{babel}
\usepackage[utf8]{inputenc}
\usepackage{pgf}
\usepackage{verbatim}
\usepackage{inconsolata}
\usepackage{listings}
\lstset{language=C, frame=single,
	basicstyle=\ttfamily,
	numbers=left, numberstyle=\tiny\color{gray},
	numbersep=5pt, fancyvrb=true
}

\usepackage[noend]{algorithm2e}
\usepackage{pgf}
\usepackage{tikz}
\usepackage{graphicx}
\usetikzlibrary{%
	arrows,
	arrows.meta,
	positioning,
	calc,
	backgrounds,
	chains,
	matrix,
	patterns,
	automata,
	fit,
	graphs,
	decorations,
	decorations.pathmorphing,
	decorations.pathreplacing,
	decorations.markings,
}

\usepgflibrary{shapes,shapes.geometric}


\usepackage{url}


\usepackage{xmpmulti}
\usepackage{transparent}
\graphicspath{{img/}}

\newcommand{\zero}{\ensuremath{\mathtt{0}}\xspace}
\newcommand{\uno}{\ensuremath{\mathtt{1}}\xspace}
\newcommand{\due}{\ensuremath{\mathtt{2}}\xspace}
\newcommand{\tre}{\ensuremath{\mathtt{3}}\xspace}
\newcommand{\kleene}{\ensuremath{{}^{\ast}}\xspace}
\xspaceaddexceptions{\zero \uno \due \tre \kleene}


\newcommand{\simplebox}[1]{
      \begin{tikzpicture}
              \node[text centered, rectangle, draw=black, thick, text width=\linewidth,align=center, rounded corners, minimum height=1em] (C) {#1};
      \end{tikzpicture}
}
%
% In questo file si mettono i comandi utili solo per questo insegnamento.
%
\author{Gianluca Della Vedova}
\title{Linguaggi e Computabilità}
\institute{Univ. Milano--Bicocca\\
	\texttt{https://www.unimib.it/gianluca-della-vedova}}
\date{\today}


\title{Linguaggi e Computabilità, Lezione 05.
  $\epsilon$-NFA}

\begin{document}


\begin{frame}
	\titlepage
\end{frame}


\tikzset{
  enfa/.style={->, >=Stealth, shorten >=1pt, auto, node distance=2.2cm,
    semithick, initial text=, every state/.style={minimum size=8mm},
    accepting/.style={double=bg, double distance=2pt}},
}



\begin{frame}
        \frametitle{Esempio 1}
	\begin{center}
		\begin{tikzpicture}[enfa]
			\node[state, initial]   (q0)               {$q_0$};
			\node[state]            (q1) [right=of q0] {$q_1$};
			\node[state]            (q2) [right=of q1] {$q_2$};
			\node[state, accepting] (q3) [right=of q2] {$q_3$};
                        \path (q0) edge                node       {$\varepsilon$}          (q1)
                                   edge [bend left=45] node       {\zero}          (q2)
			      (q1) edge [loop above]   node       {\uno}          ()
                                   edge                node       {\zero} (q2)
                              (q2) edge                node       {\uno}         (q3)
                                   edge [bend left=45] node       {\zero}          (q0);
		\end{tikzpicture}
	\end{center}
\end{frame}


\begin{frame}
        \frametitle{Esempio 2}
	\begin{center}
		\begin{tikzpicture}[enfa]
			\node[state, initial]   (q0)                     {$q_0$};
			\node[state]            (q1) [above right=1cm and 2cm of q0] {$q_1$};
			\node[state]            (q3) [below right=1cm and 2cm of q0] {$q_3$};
			\node[state, accepting] (q2) [right=of q1]       {$q_2$};
			\node[state, accepting] (q4) [right=of q3]       {$q_4$};
			\path (q0) edge node       {$\varepsilon$} (q1)
                              (q0) edge node[swap] {\uno} (q3)
                              (q1) edge node       {\uno}  (q2)
                              (q2) edge                node       {$\varepsilon$}         (q3)
                              (q3) edge node       {\zero}  (q4);
		\end{tikzpicture}
	\end{center}
\end{frame}


\begin{frame}
        \frametitle{Esempio 3}
	\begin{center}
		\begin{tikzpicture}[enfa, node distance=3cm]
			\node[state, initial, accepting] (q0)               {$q_0$};
			\node[state, accepting]          (q1) [right=of q0] {$q_1$};
			\path (q0) edge [loop above] node {$\mathtt{0}$}  ()
			           edge              node {$\varepsilon$} (q1)
			      (q1) edge [loop above] node {$\mathtt{1}$}  ();
		\end{tikzpicture}
	\end{center}
\end{frame}



\begin{frame}
        \frametitle{Esempio 4}
	\begin{center}
		\begin{tikzpicture}[enfa, node distance=3cm]
			\node[state, initial]   (q0)               {$q_0$};
			\node[state]            (q1) [right=of q0] {$q_1$};
			\node[state, accepting] (q2) [right=of q1] {$q_2$};
			\path (q0) edge              node {$\mathtt{+}, \mathtt{-}, \varepsilon$} (q1)
                              (q1) edge              node {\zero, \uno}       (q2)
                              (q2) edge [loop above] node {\zero, \uno}       ();
		\end{tikzpicture}
	\end{center}
\end{frame}

\begin{frame}
        \frametitle{Esempio 5}
	\begin{center}
		\begin{tikzpicture}[enfa]
			\node[state, initial, accepting] (q0)               {$q_0$};
			\node[state]                     (q1) [right=of q0] {$q_1$};
			\node[state]                     (q2) [right=of q1] {$q_2$};
			\node[state]                     (q3) [above right=1cm and 2cm of q2] {$q_3$};
			\node[state, accepting]          (q4) [below right=1cm and 2cm of q2] {$q_4$};
			\path (q0) edge                 node       {\uno}          (q1)
			      (q1) edge                 node       {\zero}         (q2)
			      (q2) edge [bend right=45] node[swap] {$\varepsilon$} (q0)
			           edge                 node       {\uno}          (q3)
			           edge                 node[swap] {\uno}          (q4)
			      (q3) edge [bend left=30]  node       {\zero}         (q4)
			           edge [bend right=30] node[swap] {\uno}          (q4)
			           edge [bend right=30] node[swap] {\zero}         (q2)
			      (q4) edge [bend left=30]  node       {\uno}          (q1);
		\end{tikzpicture}
	\end{center}
\end{frame}



\end{document}
